% ============================================================
\chapter{Q-Learning and Temporal Difference Methods}
\label{ch:q_learning}
% ============================================================

In 1989, Christopher Watkins, in his doctoral thesis at Cambridge, proposed a deceptively simple algorithm: at each transition, update the value of a state-action pair using the maximum over the next actions. This is \emph{Q-learning}, the first reinforcement learning algorithm proven to converge to the optimal policy \emph{without} knowing the environment model. Almost simultaneously, Rummery and Niranjan proposed \emph{SARSA}, an on-policy variant. The distinction between these two approaches---\emph{off-policy} and \emph{on-policy}---is one of the fundamental dividing lines in reinforcement learning.

\begin{intuition}
Q-Learning and SARSA are the foundational model-free control algorithms.
They learn the action-value function $Q(s,a)$ directly from experience.
SARSA is \emph{on-policy}: it evaluates the policy it follows.
Q-Learning is \emph{off-policy}: it learns the optimal policy regardless
of the exploration strategy. This chapter covers their derivation,
convergence guarantees, and exploration mechanisms.
\end{intuition}

% ------------------------------------------------------------
\section{TD Control Framework}
\label{sec:td_control}
% ------------------------------------------------------------

The \textbf{control} problem\index{control problem} is to find an optimal
policy $\pi^*$ without access to the transition model. The idea is to
combine TD estimation of $Q^\pi$ with $\epsilon$-greedy policy improvement,
in the spirit of Generalized Policy Iteration (GPI).

\begin{definition}[TD Target for Action-Values]
The one-step TD target for $Q(S_t, A_t)$ using a bootstrap estimate is:
\[
    y_t = R_{t+1} + \gamma Q(S_{t+1}, A_{t+1})
\]
where $A_{t+1}$ is the action taken in the next state. The TD error is
$\delta_t = y_t - Q(S_t, A_t)$.
\end{definition}

% ------------------------------------------------------------
\section{SARSA --- On-Policy Control}
\label{sec:sarsa}
% ------------------------------------------------------------

\begin{definition}[SARSA]
\index{SARSA}
The SARSA algorithm updates $Q(S_t, A_t)$ using the quintuplet
$(S_t, A_t, R_{t+1}, S_{t+1}, A_{t+1})$:
\[
    Q(S_t, A_t) \leftarrow Q(S_t, A_t) + \alpha \big[
    R_{t+1} + \gamma Q(S_{t+1}, A_{t+1}) - Q(S_t, A_t) \big].
\]
\end{definition}

\begin{keyformulas}[title=\textbf{SARSA Update}]
\[
    Q(S_t, A_t) \leftarrow Q(S_t, A_t) + \alpha \big[
    R_{t+1} + \gamma Q(S_{t+1}, A_{t+1}) - Q(S_t, A_t) \big]
\]
where $A_{t+1} \sim \pi_\epsilon(\cdot | S_{t+1})$ is the action actually
chosen by the $\epsilon$-greedy policy.
\end{keyformulas}

\begin{algorithme}[title=\textbf{SARSA}]
\begin{enumerate}
    \item Initialize $Q(s,a)$ arbitrarily, $Q(\text{terminal}, \cdot) = 0$
    \item For each episode:
    \begin{enumerate}
        \item Initialize $S_0$; choose $A_0$ from $\pi_\epsilon$ derived from $Q$
        \item For each step $t = 0, 1, \ldots$ until termination:
        \begin{enumerate}
            \item Execute $A_t$, observe $R_{t+1}, S_{t+1}$
            \item Choose $A_{t+1}$ from $\pi_\epsilon$ derived from $Q$
            \item $Q(S_t, A_t) \leftarrow Q(S_t, A_t) + \alpha[R_{t+1} + \gamma Q(S_{t+1}, A_{t+1}) - Q(S_t, A_t)]$
        \end{enumerate}
    \end{enumerate}
\end{enumerate}
\end{algorithme}

\begin{theorem}[SARSA Convergence]
Under the Robbins-Monro conditions and if every state-action pair $(s,a)$
is visited infinitely often (GLIE condition: Greedy in the Limit with
Infinite Exploration), SARSA converges to $Q^*$ almost surely.
\end{theorem}

% ------------------------------------------------------------
\section{Q-Learning --- Off-Policy Control}
\label{sec:q_learning}
% ------------------------------------------------------------

\begin{definition}[Q-Learning]
\index{Q-Learning}
The Q-Learning algorithm (Watkins, 1989) updates:
\[
    Q(S_t, A_t) \leftarrow Q(S_t, A_t) + \alpha \big[
    R_{t+1} + \gamma \max_{a'} Q(S_{t+1}, a') - Q(S_t, A_t) \big].
\]
\end{definition}

\begin{keyformulas}[title=\textbf{Q-Learning Update}]
\[
    Q(S_t, A_t) \leftarrow Q(S_t, A_t) + \alpha \Big[
    R_{t+1} + \gamma \max_{a' \in \mathcal{A}} Q(S_{t+1}, a') - Q(S_t, A_t) \Big]
\]
\end{keyformulas}

The key difference with SARSA is the use of $\max$ instead of
$Q(S_{t+1}, A_{t+1})$. Q-Learning learns $Q^*$ directly, regardless of
the behavior policy (as long as all pairs $(s,a)$ are visited).

\begin{algorithme}[title=\textbf{Q-Learning}]
\begin{enumerate}
    \item Initialize $Q(s,a)$ arbitrarily, $Q(\text{terminal}, \cdot) = 0$
    \item For each episode:
    \begin{enumerate}
        \item Initialize $S_0$
        \item For each step $t = 0, 1, \ldots$ until termination:
        \begin{enumerate}
            \item Choose $A_t$ from $\pi_\epsilon$ derived from $Q$
            \item Execute $A_t$, observe $R_{t+1}, S_{t+1}$
            \item $Q(S_t, A_t) \leftarrow Q(S_t, A_t) + \alpha[R_{t+1} + \gamma \max_{a'} Q(S_{t+1}, a') - Q(S_t, A_t)]$
        \end{enumerate}
    \end{enumerate}
\end{enumerate}
\end{algorithme}

\begin{theorem}[Q-Learning Convergence]
\label{thm:q_learning_conv}
If every pair $(s,a)$ is visited infinitely often and the learning rates
$\alpha_t(s,a)$ satisfy the Robbins-Monro conditions:
\[
    \sum_t \alpha_t(s,a) = \infty, \qquad \sum_t \alpha_t^2(s,a) < \infty,
\]
then $Q(s,a) \to Q^*(s,a)$ almost surely for all $(s,a)$.
\end{theorem}

\begin{proof}[Proof sketch]
The update is a stochastic approximation of the Bellman optimality
operator $\mathcal{T}^*$. The contraction property of $\mathcal{T}^*$
in the $\ell_\infty$ norm, combined with the Robbins-Monro conditions,
guarantees convergence (Jaakkola et al., 1994).
\end{proof}

% ------------------------------------------------------------
\section{SARSA vs Q-Learning: The Cliff Walking Example}
\label{sec:cliff}
% ------------------------------------------------------------

\begin{example}[Cliff Walking]
\index{cliff walking}
The \texttt{CliffWalking-v0} environment illustrates the difference:
\begin{itemize}
    \item \textbf{Q-Learning} learns the optimal policy (along the cliff edge),
    but the $\epsilon$-greedy exploration may cause the agent to fall.
    \item \textbf{SARSA} learns a safer policy (away from the cliff) because
    it accounts for the $\epsilon$-greedy behavior in its updates.
\end{itemize}
Q-Learning optimizes the return of the \emph{target} (greedy) policy, while
SARSA optimizes the return of the \emph{actual} ($\epsilon$-greedy) policy.
\end{example}

% ------------------------------------------------------------
\section{Expected SARSA}
\label{sec:expected_sarsa}
% ------------------------------------------------------------

\begin{definition}[Expected SARSA]
\index{Expected SARSA}
\textbf{Expected SARSA} uses the expectation over next actions instead
of a sampled action:
\[
    Q(S_t, A_t) \leftarrow Q(S_t, A_t) + \alpha \Big[
    R_{t+1} + \gamma \sum_{a'} \pi(a' | S_{t+1}) Q(S_{t+1}, a') - Q(S_t, A_t) \Big].
\]
\end{definition}

\begin{remark}
Expected SARSA generalizes both SARSA and Q-Learning:
\begin{itemize}
    \item If $\pi$ is the current $\epsilon$-greedy policy, it is an improved
    version of SARSA with lower variance.
    \item If $\pi$ is the greedy policy, it reduces to Q-Learning.
\end{itemize}
\end{remark}

% ------------------------------------------------------------
\section{$n$-Step Methods}
\label{sec:nstep}
% ------------------------------------------------------------

\begin{definition}[$n$-Step Return]
The $n$-step return combines $n$ immediate rewards with a bootstrap estimate:
\[
    G_t^{(n)} = \sum_{k=0}^{n-1} \gamma^k R_{t+k+1} + \gamma^n Q(S_{t+n}, A_{t+n}).
\]
\end{definition}

\begin{proposition}[Bias-Variance of $n$-Step Returns]
\begin{itemize}
    \item $n=1$: TD(0) --- low variance, high bias.
    \item $n=\infty$: Monte Carlo --- zero bias, high variance.
    \item Intermediate $n$: trade-off between bias and variance.
\end{itemize}
In practice, $n \in \{4, 8, 16\}$ often works well.
\end{proposition}

% ------------------------------------------------------------
\section{Exploration Strategies}
\label{sec:exploration}
% ------------------------------------------------------------

\begin{definition}[$\epsilon$-Greedy Exploration]
\index{epsilon-greedy}
The $\epsilon$-greedy policy derived from $Q$ is:
\[
    \pi_\epsilon(a \mid s) = \begin{cases}
        1 - \epsilon + \frac{\epsilon}{\abs{\mathcal{A}}} & \text{if } a = \argmax_{a'} Q(s,a') \\
        \frac{\epsilon}{\abs{\mathcal{A}}} & \text{otherwise}
    \end{cases}
\]
Typically $\epsilon$ is decayed over time: $\epsilon_t = \max(\epsilon_{\min}, \epsilon_0 \cdot \beta^t)$.
\end{definition}

\begin{definition}[Boltzmann (Softmax) Exploration]
\index{Boltzmann exploration}
Boltzmann exploration selects actions with probability proportional to
exponentiated Q-values:
\[
    \pi_\tau(a \mid s) = \frac{\exp(Q(s,a) / \tau)}{\sum_{a'} \exp(Q(s,a') / \tau)}
\]
where $\tau > 0$ is the temperature. As $\tau \to 0$, this becomes greedy;
as $\tau \to \infty$, it becomes uniform.
\end{definition}

\begin{definition}[Upper Confidence Bound (UCB)]
\index{UCB}
UCB-based exploration selects:
\[
    A_t = \argmax_a \left[ Q(s,a) + c \sqrt{\frac{\ln N(s)}{N(s,a)}} \right]
\]
where $N(s)$ and $N(s,a)$ are visit counts and $c > 0$ controls the
exploration bonus.
\end{definition}

\begin{attention}[title=\textbf{Exploration-Exploitation Dilemma}]
Insufficient exploration leads to sub-optimal convergence (the agent may
miss high-reward regions). Excessive exploration wastes samples on
uninformative actions. The optimal balance depends on the environment
structure and the learning horizon.
\end{attention}

% ------------------------------------------------------------
\section{Maximization Bias and Double Q-Learning}
\label{sec:double_q}
% ------------------------------------------------------------

\begin{definition}[Maximization Bias]
\index{maximization bias}
Q-Learning uses $\max_{a'} Q(S_{t+1}, a')$ as the target. Because $Q$
is a noisy estimate, the $\max$ introduces a positive bias:
$\E[\max_a Q(s,a)] \geq \max_a \E[Q(s,a)]$.
\end{definition}

\begin{definition}[Double Q-Learning]
\index{Double Q-Learning}
Double Q-Learning (Hasselt, 2010) maintains two independent estimators
$Q_1, Q_2$. At each step, one is randomly chosen for the update:
\[
    Q_1(S_t, A_t) \leftarrow Q_1(S_t, A_t) + \alpha \Big[
    R_{t+1} + \gamma Q_2\big(S_{t+1}, \argmax_{a'} Q_1(S_{t+1}, a')\big) - Q_1(S_t, A_t)
    \Big].
\]
The roles of $Q_1$ and $Q_2$ are swapped with probability $0.5$.
\end{definition}

\begin{proposition}[Double Q-Learning removes maximization bias]
By decoupling action selection ($\argmax$ on $Q_1$) from evaluation
(using $Q_2$), Double Q-Learning produces unbiased estimates in the limit,
eliminating the systematic overestimation of standard Q-Learning.
\end{proposition}

% ------------------------------------------------------------
\section{Python Implementation}
\label{sec:q_impl}
% ------------------------------------------------------------

\begin{codeblock}[title=\textbf{Q-Learning Implementation}]
\begin{minted}{python}
import numpy as np
import gymnasium as gym

def q_learning(env, n_episodes=10000, alpha=0.1,
               gamma=0.99, epsilon=1.0, eps_decay=0.999,
               eps_min=0.01):
    """Tabular Q-learning with epsilon-greedy exploration."""
    nS = env.observation_space.n
    nA = env.action_space.n
    Q = np.zeros((nS, nA))

    for ep in range(n_episodes):
        s, _ = env.reset()
        done = False
        while not done:
            # epsilon-greedy action selection
            if np.random.random() < epsilon:
                a = env.action_space.sample()
            else:
                a = np.argmax(Q[s])
            s_next, r, terminated, truncated, _ = env.step(a)
            done = terminated or truncated
            Q_next = 0.0 if terminated else np.max(Q[s_next])
            Q[s, a] += alpha * (r + gamma * Q_next - Q[s, a])
            s = s_next
        epsilon = max(eps_min, epsilon * eps_decay)
    return Q
\end{minted}
\end{codeblock}

\begin{codeblock}[title=\textbf{SARSA Implementation}]
\begin{minted}{python}
def sarsa(env, n_episodes=10000, alpha=0.1,
          gamma=0.99, epsilon=0.1):
    """Tabular SARSA with epsilon-greedy policy."""
    nS = env.observation_space.n
    nA = env.action_space.n
    Q = np.zeros((nS, nA))

    def eps_greedy(s):
        if np.random.random() < epsilon:
            return env.action_space.sample()
        return np.argmax(Q[s])

    for ep in range(n_episodes):
        s, _ = env.reset()
        a = eps_greedy(s)
        done = False
        while not done:
            s_next, r, terminated, truncated, _ = env.step(a)
            done = terminated or truncated
            a_next = eps_greedy(s_next) if not done else 0
            Q_next = 0.0 if terminated else Q[s_next, a_next]
            Q[s, a] += alpha * (r + gamma * Q_next - Q[s, a])
            s, a = s_next, a_next
    return Q
\end{minted}
\end{codeblock}

% ------------------------------------------------------------
\section{Exercises}
\label{sec:q_exercises}
% ------------------------------------------------------------

\begin{exercise}[Cliff Walking Comparison]
Implement Q-Learning and SARSA on \texttt{CliffWalking-v0}. Plot the
per-episode reward for both algorithms. Explain the difference in learned
policies.
\end{exercise}

\begin{exercise}[Exploration Strategies]
Compare $\epsilon$-greedy, Boltzmann, and UCB exploration on
\texttt{FrozenLake-v1}. Measure convergence speed (episodes to reach
$>0.7$ average reward) for each strategy.
\end{exercise}

\begin{exercise}[Double Q-Learning]
Implement Double Q-Learning and compare with standard Q-Learning on an
MDP where maximization bias is harmful (e.g., the ``Maximization Bias
Example'' from Sutton \& Barto, Chapter 6). Plot the percentage of times
the correct action is chosen at the start state.
\end{exercise}

\begin{exercise}[Expected SARSA Derivation]
Show that Expected SARSA with a greedy target policy is equivalent to
Q-Learning. What happens when $\epsilon = 0$ in the $\epsilon$-greedy case?
\end{exercise}
